\section{Benchmarking Real-World Agents}

\subsection{Benchmark Pipeline}

Anthropic 把 Agent 能力拆解成多个指标：自主性、工具调用、推理深度、偏好一致性。每轮评估都会串联以下环节：

\begin{enumerate}
    \item \textbf{Design scenarios}：选定长周期任务（调研、计划、执行、汇报）
    \item \textbf{Establish baselines}：用人类团队或前代模型做对照
    \item \textbf{Instrument data}：收集 agent 的行动日志、API 调用、审计 trail
    \item \textbf{Score metrics}：自动化评分 + 人类回顾
    \item \textbf{Feedback to training}：将发现的 failure modes 逐条转化为 dataset/behavioral trace
\end{enumerate}

\begin{figure}[H]
\centering
\includegraphics[page=3,width=0.92\textwidth]{slides.pdf}
\caption{Agent Benchmark pipeline illustrated on slide 3.\protect\footnotemark}
\end{figure}
\footnotetext{Slides emphasize multi-step evaluation loops.}

\subsection{Case Study: Decision Support Task}

一个典型任务是让 Agent 读完 20 页内部战略文档，生成 POAP（Plan, Options, Actions, Priority）摘要并安排 follow-up 会议。

\begin{itemize}
    \item Agent 需要跨页面保持 coherence（段落嵌套、引用）
    \item 输出同时要兼顾事实、推理链、行动建议
    \item 最终由人类 reviewer 根据 `POAP` template 打分
\end{itemize}

\begin{knowledgebox}{Tailored judgment rubric}
Anthropic 为每种任务制定 rubric：事实准确性、结构完整性、工具调用合规性、是否提供可操作建议。只有在所有维度都合格的情况下才会推进到 ASL 升级。
\end{knowledgebox}

\subsection{本章小结}

Anthropic 的 benchmark pipeline 把多个任务场景、数据采集、评分和训练反馈串成一个闭环；每个 scenario 都配套专门的 rubric，从而把能力评估变成可操作的治理信号。

\subsection{Evidence Matrix}

为了把能力评估转成决策信息，Anthropic 使用\textbf{evidence matrix}：按任务阶段列出关键指标、数据源、当前表现，并给出安全边界（green/orange/red）。

\begin{table}[H]
\centering
\begin{tabular}{p{0.18\textwidth}p{0.24\textwidth}p{0.24\textwidth}p{0.24\textwidth}}
\toprule
阶段 & 指标 & 数据源 & 当前风险信号 \\
\midrule
Planning & clarity of plan, dependency awareness & prompt logs, human evaluation & plan misses key stakeholders \\
Tool invocation & tool success rate, error rate & API logs, retry counts & spike in `api\_error` events \\
Execution & prospect completion quality & manager reviews, outcome diff & 5\% of tasks reworked \\
Reporting & factual accuracy & summary diff + external documents & hallucinated reference detected \\
\bottomrule
\end{tabular}
\caption{Evidence matrix used for capability snapshots}
\end{table}

\begin{warningbox}{Evidence matrix refresh cadence}
Matrix rows are refreshed every week or after any significant dataset/model update. Without timely refresh, the safety gates rely on stale performance snapshots.
\end{warningbox}

